\section{Conclusiones y Trabajos Futuros}

\subsection{6.1 Conclusiones}

Resulta revelador observar cómo un problema que hace apenas una década parecía lejano se ha convertido en uno de los principales cuellos de botella para la sostenibilidad de las operaciones en LEO. Este trabajo ha presentado una arquitectura híbrida ground-onboard orientada a elevar significativamente los estándares de seguridad orbital en mega-constelaciones de telecomunicaciones.

Los resultados de simulación demuestran que la integración de capacidades autónomas, apoyadas en star trackers repurposed, permite reducciones sustanciales tanto en probabilidad de colisión residual como en consumo de propelente. La incorporación de detección óptica embarcada, tal como validan estudios recientes \cite{campiti2025detectability}, aporta una capa de resiliencia que los sistemas puramente ground-based difícilmente pueden igualar en escenarios de alta densidad o bajo enlace de comunicaciones.

Más allá de las métricas numéricas, la propuesta contribuye a un cambio de paradigma: pasar de una gestión reactiva y centralizada hacia un modelo híbrido más distribuido y adaptable. Esta transición no elimina la necesidad de coordinación global, pero la complementa de forma inteligente, ofreciendo mayor robustez operativa. Después de años siguiendo el debate sobre congestión orbital, considero que avances de este tipo son necesarios —aunque no suficientes por sí solos— para que el crecimiento de las constelaciones de telecomunicaciones sea compatible con la preservación del entorno espacial a largo plazo.

\subsection{6.2 Trabajos Futuros}

Lejos de considerar el trabajo completado, varias líneas abiertas merecen atención prioritaria. Particularmente prometedora resulta la extensión del marco de aprendizaje por refuerzo hacia entornos multi-agente, donde cientos o miles de satélites coordinen maniobras de forma distribuida minimizando interferencias mutuas. Investigaciones recientes en planificación autónoma mediante reinforcement learning y on-orbit servicing \cite{patnala2025oos} ofrecen una base sólida sobre la cual construir algoritmos que consideren no solo la evasión individual, sino la estabilidad colectiva de la constelación.

Entre las prioridades identificadas destacan: la validación en vuelo de los módulos onboard en misiones demostradoras, la integración con sistemas de rendezvous y servicing para maniobras más complejas con bajo consumo de propelente, y el desarrollo de protocolos estandarizados de transparencia para agentes autónomos. También resulta crítico explorar el comportamiento del sistema bajo condiciones extremas de actividad solar y el impacto de degradación gradual de sensores a lo largo de la vida útil del satélite.

Finalmente, creemos necesario avanzar en el diseño de benchmarks comunes que permitan comparar diferentes aproximaciones de evasión autónoma bajo escenarios estandarizados. Solo mediante este esfuerzo colectivo —técnico y regulatorio— lograremos que la gestión del tráfico espacial evolucione al ritmo que demanda el rápido crecimiento de actividades en órbita baja.